% Standalone popularising summary for the master's thesis.
% Compiles on its own with pdflatex. It is deliberately kept separate from the
% main thesis so it can be submitted and read as its own short document.
\documentclass[11pt]{article}

\usepackage[a4paper,margin=2.5cm]{geometry}
\usepackage[T1]{fontenc}
\usepackage[utf8]{inputenc}
\usepackage{lmodern}
\usepackage{microtype}
\usepackage{parskip}

\title{Popularising Summary}
\author{Louis Vandenbruwaene}
\date{}

\begin{document}

% The faculty caps this summary at one A4 page and 3,500 characters, so the
% title block is compact and the text is trimmed to roughly 3,450 characters.
{\large\bfseries Popularising Summary}\\[2pt]
{\normalsize Louis Vandenbruwaene}\\[6pt]
\textit{Master's thesis: an extremal problem on connectivity in networks (Erd\H{o}s Problem 915).}

\bigskip

Many of the networks we rely on are useful precisely because they offer backup routes. If one road is closed, traffic takes another. If one cable fails, a message still arrives along a different chain of links. Mathematicians study this with graphs, where dots stand for objects and lines for the connections between them. This thesis asks a simple-sounding question: how many connections can a network with a fixed number of points have, if no two of those points may be joined by too many separate, non-overlapping routes?

That question is a version of a problem posed decades ago by B\'{e}la Bollob\'{a}s and Paul Erd\H{o}s, listed today as Erd\H{o}s Problem 915. In its original form it concerns ordinary networks and routes that may not share a connection, and for that case the answer has been known since a theorem of Mader. This thesis starts from that result and then changes the rules of the network one at a time: connections may be doubled up, directions may be added so that a link works only one way, or a single connection may join more than two points at once. The variants need different proofs and computational checks.

At that point the thesis hands the work to a computer. One program describes every version of the network in the same way, measures exactly how many separate routes run between any two points, and proves the answer outright for small networks by checking every case. Where checking everything is hopeless, it searches for good examples instead, walking from one network to a slightly better one while keeping a short memory of its recent moves so it cannot simply undo them. When the search is pointed at cases we already understand, it rediscovers the very networks the earlier proofs singled out.

The interplay between proof and experiment then pays off in five places, and the thesis is explicit about how far each one is checked. Arguments drafted with AI assistance that the author has not verified line by line are presented as conjectures with a proposed proof attached, rather than as settled theorems. For networks whose connections may each tie several points together at once, the strictest limit on backup routes is settled exactly at every size, and for the next limit the thesis proposes a formula as an upper limit, attained in a stated range. For one-way networks, recent work of others already fixes the largest number of links at a quarter of the square of the number of points, give or take an amount growing only in proportion to that number, and the thesis adds a self-contained argument, not yet verified, that also gives a bound at every size. It exhibits dense examples without asserting their optimality, and a counterexample rules out an earlier formula for all sizes. For networks whose links may be doubled up, and whose routes must avoid sharing any intermediate point as well, the thesis poses a question that had not been separated out before and settles it exactly under the strictest limit. Beyond that it proposes that the answer is decided by the best densely connected piece of each size. For one-way networks whose links may be doubled up, it proposes an exact answer at every size. The idea is to find, inside any such network, a thin skeleton that alone remembers who can reach whom, and to peel off one skeleton for each unit of the limit. And for one-way networks with many-point connections, a two-layer design is proposed to have the optimal leading rate of growth.

\end{document}
